% ==============================================================================
% ================================ Zusammenfassung ==============================
% ==============================================================================

\chapter{Zusammenfassung und Reflexion}
\label{chap:zusammenfassung}
\thispagestyle{fancy}

\section{Reflexion}

Die Entwicklung zeigte, dass die Kombination aus \acs{PHP} und einer relationalen Datenbank für einfache Webanwendungen gut geeignet ist. Besonders hilfreich war die Verwendung von \acs{PDO} für eine einheitliche und sichere Datenbankanbindung.

Herausforderungen ergaben sich bei der Validierung von Nutzereingaben und der Unterscheidung zwischen anonymen und registrierten Teilnehmern. Die Export-Funktion von phpMyAdmin erwies sich als praktikable Lösung zur Sicherung des Datenbankzustands.

Die Arbeit hat gezeigt, dass eine strukturierte Vorgehensweise – von der Datenbankmodellierung über die Implementierung bis zur Dokumentation – für die Entwicklung funktionaler Anwendungen essenziell ist. Die gewählte Architektur ist wartbar und kann um weitere Funktionen erweitert werden:

\begin{forest}
    for tree={
        font=\ttfamily,
        grow'=0,
        child anchor=west,
        parent anchor=south,
        anchor=west,
        calign=first,
        edge path={
            \noexpand\path [draw, \forestoption{edge}] (!u.south west) +(10pt,0) |- node[fill,inner sep=1.5pt] {} (.child anchor)\forestoption{edge label};
        },
        before typesetting nodes={
            if n=1{
                insert before={[,phantom,calign with current]}
            }{}
        },
        fit=band,
        before computing xy={l=22pt},
    }
    [\textbf{C:\textbackslash xampp\textbackslash htdocs\textbackslash survey}
        [\textbf{index.php}]
        [\textbf{report.php}]
        [\textbf{save.php}]
        [\textbf{survey.php}]
        [\textbf{inc}
                [\textbf{db.php}]
                [\textbf{db.sql}]
        ]
    ]
\end{forest}

% ==============================================================================
% ================================ Zusammenfassung ==============================
% ==============================================================================